% !TEX TS-program = xelatex
% !TEX encoding = UTF-8 Unicode
% !Mode:: "TeX:UTF-8"

\documentclass{resume}
\usepackage{zh_CN-Adobefonts_external} % Simplified Chinese Support using external fonts (./fonts/zh_CN-Adobe/)
%\usepackage{zh_CN-Adobefonts_internal} % Simplified Chinese Support using system fonts
\usepackage{linespacing_fix} % disable extra space before next section
\usepackage{cite}
\usepackage[hidelinks]{hyperref}

\begin{document}
\pagenumbering{gobble} % suppress displaying page number

\name{秦浩翔}

% {E-mail}{mobilephone}{homepage}
% be careful of _ in emaill address
\contactInfo{(+86) 181-0805-1863}{qhx\_sjtu@sjtu.edu.cn}{24级软件工程本科生}{GitHub \href{https://github.com/1JI0O}{@1JI0O}}

% {E-mail}{mobilephone}
% keep the last empty braces!
%\contactInfo{xxx@yuanbin.me}{(+86) 131-221-87xxx}{}
 
\section{个人总结}
我有扎实的编程基础，熟悉python和c++，以及能够熟练使用codex和claude code等code agent，此外我也具备丰富的linux使用经验，以及一定的openclaw使用经验。我对多模态，计算机视觉，具身智能等领域有浓厚兴趣，曾进行过相关课程的自学和多个项目实践，有不错的自学和动手实操能力。我目前的未来规划是出国深造，但是我最终还是希望进入工业界工作

\vspace{0.5ex}

\href{https://1ji0o.github.io/}{我的博客} 收录了我的部分课程学习笔记和项目工程日志等内容

% \section{\faGraduationCap\ 教育背景}
\section{教育背景}
% \datedsubsection{\textbf{中国科学院大学},计算机应用技术,\textit{在读硕士研究生}}{2015.9 - 2018.6}
% \ \textbf{排名11/133(前10\%)},中国科学院大学学业奖学金(2次),IEEE Student member,预计2018年6月毕业
\datedsubsection{\textbf{上海交通大学},软件工程,\textit{本科}}{2024.9 - 现在}
\ 核心课程学积分 85.639470199	   排名 44 / 93

\ 2024年度致远荣誉奖学金


% % \section{\faCogs\ IT 技能}
% \section{技术能力}
% % increase linespacing [parsep=0.5ex]
% \begin{itemize}[parsep=0.2ex]
%   \item \textbf{编程语言}: JavaScript (ECMAScript, Node.js), HTML/CSS, Python, Go, SQL, C, Shell
%   \item \textbf{操作系统,数据库与工程构建}: Linux/macOS/MySQL/MongoDB/Git/webpack/Progressive Web App
%   \item \textbf{关键词}: React/Vue.js/D3.js(SVG)/three.js(canvas, WebGL)/chrome extension/Express
% \end{itemize}

% \end{itemize}

\section{实习经历}
\datedsubsection{\textbf{MVIG实验室}, 本科实习生}{2025.12 - 2026.5}
\begin{itemize}
  \item 承担了Force Policy (RSS,2026)的部分实验任务(跑实验和适配了一个baseline)
  \item 在学长指导下独立开展了对一个想法的研究
\end{itemize}


% \begin{onehalfspacing}
% \end{onehalfspacing}

% \datedsubsection{\textbf{DID-ACTE} 荷兰莱顿}{2015年}
% \role{本科毕业设计}{LIACS 交换生}
% 利用结巴分词对中国古文进行分词与词性标注，用已有领域知识训练形成 classifier 并对结果进行调优
% \begin{onehalfspacing}
% \begin{itemize}
%   \item 利用结巴分词对中国古文进行分词与词性标注
%   \item 利用已有领域知识训练形成 classifier, 并用分词结果进行测试反馈
%   \item 尝试不同规则，对 classifier 进行调优
% \end{itemize}
% \end{onehalfspacing}

\section{论文/项目作品}
\textbf{Force Policy: Learning Hybrid Force-Position Control Policy under Interaction Frame for Contact-Rich Manipulation}

Hongjie Fang*, Shirun Tang*, Mingyu Mei*, \textbf{Haoxiang Qin}, Zihao He, Jingjing Chen, Ying Feng, Chenxi Wang, Wanxi Liu, Zaixing He, Cewu Lu, Shiquan Wang$^\dagger$

RSS,2026 \href{https://arxiv.org/abs/2602.22088}{[Paper]} \href{https://force-policy.github.io/}{[Project]}

\vspace{1.0ex}

\textbf{对一个想法的研究}

核心是让动作策略模型忽略相机观测中机械臂部分，避免模型过拟合到机械臂上。研究中使用了SAM2、SAM3和ProPainter等项目，并对SAM2进行了微调，以及对原始动作策略模型\href{https://github.com/1JI0O/RISE-2/}{RISE-2}的ViT等部分进行了修改

\vspace{1.0ex}

\textbf{软件学院SEP课程大作业}

QLink和QBasic : 基于Qt的连连看游戏和gui Basic解释器。其中QLink我分别用c++和python实现了两版。这两个项目我都开源放在github上了

\vspace{0.5ex}

% \section{\faHeartO\ 项目/作品摘要}
% \section{项目/作品摘要}
% \datedline{\textit{An Integrated Version of Security Monitor Vis System}, https://hijiangtao.github.io/ss-vis-component/ }{}
% \datedline{\textit{Dark-Tech}, https://github.com/hijiangtao/dark-tech/ }{}
% \datedline{\textit{融合社交网络数据挖掘的电视节目可视化分析系统}, https://hijiangtao.github.io/variety-show-hot-spot-vis/}{}
% \datedline{\textit{LeetCodeOJ Solutions}, https://github.com/hijiangtao/LeetCodeOJ}{}
% \datedline{\textit{Info-Vis}, https://github.com/ISCAS-VIS/infovis-ucas}{}


% \section{\faInfo\ 社会实践/其他}
\section{课程学习}
% increase linespacing [parsep=0.5ex]
% \begin{itemize}[parsep=0.2ex]
%   \item 乐于参与开源社区讨论,\textbf{参与翻译 Vue.js, webpack, WebAssembly, Babel 文档，印记中文成员}
%   \item 中国科学院大学2016秋季学期可视化与可视分析课程助教，\textit{http://vis.ios.ac.cn/infovis-ucas/}
%   \item 未来论坛学生会成员、北理社联新闻信息中心主任、北理工软件学院学生会宣传部副部长(2012-2016)
%   \item 2013-2015 北京市共青团“温暖衣冬”志愿者,第九届园博会志愿者,2014 FLL机器人世锦赛志愿者
% \end{itemize}
\textbf{CS231n}

我在2024.10-2024.11期间自学了2017版CS231n课程，完成了2025年课程作业


%% Reference
%\newpage
%\bibliographystyle{IEEETran}
%\bibliography{mycite}
\end{document}
